% theme: aqueous_solutions_and_ions
\MajorQuestion{水溶液とイオン}
\SetRowSpaceQanda

\Instruction{原子のつくりに関する次の問いに答えなさい。}
\QQ{原子は、原子核とそのまわりを動く何からできているか。}{電子}
\QQ{原子核は、陽子と何からできているか。}{中性子}
\QQ{陽子がもつ電気は、正と負のどちらか。}{正の電気}
\QQ{電子がもつ電気は、正と負のどちらか。}{負の電気}
\QQ{同じ元素で、中性子の数が異なる原子どうしを何というか。}{同位体}

\Instruction{イオンに関する次の問いに答えなさい。}
\QQ{原子が電気を帯びたものを\NamedBlank{ion}という。}{イオン}
\QQ{原子が電子を失い、正の電気を帯びたものを何というか。}{陽イオン}
\QQ{原子が電子を受けとり、負の電気を帯びたものを何というか。}{陰イオン}
\QQ{水素イオンを、イオンの記号で表しなさい。}{$\mathrm{H^+}$}
\QQ{塩化物イオンを、イオンの記号で表しなさい。}{$\mathrm{Cl^-}$}

\Instruction{水溶液と電離に関する次の問いに答えなさい。}
\QQ{水に溶けたとき、水溶液に電流が流れる物質を\NamedBlank{electrolyte}という。}{電解質}
\QQ{水に溶けても、水溶液に電流が流れない物質を何というか。}{非電解質}
\QQ{物質が水に溶けて、陽イオンと陰イオンに分かれることを\NamedBlank{ionization}という。}{電離}
\QQ{塩化ナトリウムの\RefBlank{ionization}を、化学式で表しなさい。}{$\mathrm{NaCl \rightarrow Na^+ + Cl^-}$}
\QQ{砂糖とエタノールは、\RefBlank{electrolyte}と非電解質のどちらか。}{非電解質}

\Instruction{電気分解に関する次の問いに答えなさい。}
\QQ{塩化銅水溶液に電流を流すと、陰極に付着する物質は何か。}{銅}
\QQ{塩化銅水溶液に電流を流すと、陽極から発生する気体は何か。}{塩素}
\QQ{塩化銅の\RefBlank{ionization}を、化学式で表しなさい。}{$\mathrm{CuCl_2 \rightarrow Cu^{2+} + 2Cl^-}$}
\QQ{塩酸に電流を流すと、陰極から発生する気体は何か。}{水素}
\QQ{塩酸に電流を流すと、陽極から発生する気体は何か。}{塩素}
